\documentclass[11pt,a4paper]{article}
\usepackage[utf8]{inputenc}
\usepackage[T1]{fontenc}
\usepackage[catalan]{babel}
\usepackage{lmodern}
\usepackage[sfdefault,lining]{FiraSans}
\usepackage[a4paper,top=2.3cm,bottom=2.2cm,left=1.8cm,right=1.8cm,headheight=1.3cm,headsep=0.4cm,footskip=1cm]{geometry}
\usepackage{amsmath,amssymb}
\usepackage{icomma}
\usepackage{graphicx}
\usepackage[table]{xcolor}
\usepackage{tikz}
\usepackage{fontawesome5}
\usepackage{ccicons}
\usepackage{fancyhdr}
\usepackage{lastpage}
\usepackage{tasks}
\usepackage{needspace}
\usepackage[most]{tcolorbox}
\usepackage{array}
\usepackage[hidelinks,pdftitle={Fitxa d'exercicis de la unitat 10: Estadística. 1r d'ESO},pdfsubject={Matemàtiques 1r ESO},pdfkeywords={estadística, població, mostra, variables estadístiques, freqüències, taules de freqüències, diagrama de barres, polígon de freqüències, histograma, diagrama de sectors, pictograma, moda, mitjana, mediana, rang, desviació mitjana}]{hyperref}

\definecolor{navy}{HTML}{1B2748}
\definecolor{accent}{HTML}{E8603A}
\definecolor{teal}{HTML}{1F9E90}
\definecolor{soft}{HTML}{EEF1F8}
\definecolor{muted}{HTML}{5D6685}

\setlength{\parindent}{0pt}
\setlength{\parskip}{3pt}
\settasks{label=\alph*), label-format=\bfseries\color{navy}, label-width=1.3em, item-indent=1.8em, column-sep=0.6em, after-item-skip=5pt, before-skip=3pt, after-skip=4pt}

% ---------- capçalera i peu ----------
\newcommand{\llicencia}{\href{https://creativecommons.org/licenses/by-nc-sa/4.0/deed.ca}{\ccbyncsa}\ \ Llicència Creative Commons Reconeixement-NoComercial-CompartirIgual 4.0 Internacional (CC BY-NC-SA 4.0)}
\pagestyle{fancy}
\fancyhf{}
\fancyhead[L]{\raisebox{-0.35\height}{\includegraphics[height=0.95cm]{logo-gva.png}}}
\fancyhead[C]{\small\color{muted}Fitxa d'exercicis · Unitat 10 · 1r d'ESO}
\fancyhead[R]{\raisebox{-0.35\height}{\includegraphics[height=0.95cm]{logo-vermellar.png}}}
\fancyfoot[L]{\scriptsize\color{muted}\llicencia}
\fancyfoot[R]{\small\color{muted}\thepage}
\renewcommand{\headrulewidth}{0.4pt}
\renewcommand{\headrule}{\color{muted!40}\hrule height \headrulewidth}
\fancypagestyle{primera}{\fancyhf{}\renewcommand{\headrulewidth}{0pt}
  \fancyfoot[L]{\scriptsize\color{muted}\llicencia}
  \fancyfoot[R]{\small\color{muted}\thepage}}

% ---------- nivells i exercicis ----------
\newcommand{\estrelles}[1]{{\color{accent}\ifcase#1\or\faStar\faStar[regular]\faStar[regular]\or\faStar\faStar\faStar[regular]\or\faStar\faStar\faStar\fi}}
\newcounter{exer}
\newcommand{\ex}[1]{\par\needspace{4\baselineskip}\medskip\stepcounter{exer}\noindent{\color{navy}\bfseries\theexer.}\ \ {\small\estrelles{#1}}\ \ }

\newcommand{\seccio}[2]{\par\needspace{10\baselineskip}\bigskip
  \begin{tcolorbox}[colback=navy,colframe=navy,coltext=white,arc=2mm,boxrule=0pt,left=3mm,right=3mm,top=1.2mm,bottom=1.2mm]
  {\large\bfseries #1\quad #2}
  \end{tcolorbox}\vspace{-2pt}}

\newcolumntype{C}[1]{>{\centering\arraybackslash}p{#1}}
\definecolor{violet}{HTML}{7C6CF0}
\definecolor{gold}{HTML}{E0A21B}

% quadre de recordatori (només a la fitxa d'exercicis)
\newcommand{\recorda}[1]{\par\smallskip\begin{tcolorbox}[colback=accent!7,colframe=accent!55,arc=2mm,boxrule=0.5pt,left=2.5mm,right=2.5mm,top=1mm,bottom=1mm]\small{\color{accent}\faLightbulb[regular]}\ \textbf{Recorda:}\ #1\end{tcolorbox}}
\newcommand{\res}[1]{#1}\newcommand{\rest}[1]{#1}

\def\fmneed{14}
% =====================================================================
% Capçalera FentMates de les fitxes: un color per a cada curs.
% Els colors són els de la pàgina inicial (LEVELS a assets/js/materials.js);
% el logo ve de _eines/logo/pdf/ (python3 _eines/logo/genera.py).
%
% Ús, al final del preàmbul de la unitat:
%   % =====================================================================
% Capçalera FentMates de les fitxes: un color per a cada curs.
% Els colors són els de la pàgina inicial (LEVELS a assets/js/materials.js);
% el logo ve de _eines/logo/pdf/ (python3 _eines/logo/genera.py).
%
% Ús, al final del preàmbul de la unitat:
%   % =====================================================================
% Capçalera FentMates de les fitxes: un color per a cada curs.
% Els colors són els de la pàgina inicial (LEVELS a assets/js/materials.js);
% el logo ve de _eines/logo/pdf/ (python3 _eines/logo/genera.py).
%
% Ús, al final del preàmbul de la unitat:
%   \input{../comu/fentmates.tex}
%   \fmcurs{1eso}
% i, al començament del document:
%   \capcalera{Matemàtiques}{1r ESO}{Fitxa d'exercicis · Unitat 2}{Divisibilitat}
% Opcional abans de \input: \def\fmneed{14} (espai mínim abans de cada secció).
% =====================================================================
\makeatletter
\def\fm@logodir{../../logo/pdf}
\providecommand{\fmneed}{10}
\def\fm@cursnom{}

% clau del curs, color i color del text damunt del color
\newcommand{\fm@curs}[3]{\@namedef{fm@color@#1}{#2}\@namedef{fm@text@#1}{#3}}
\fm@curs{1eso}{E6B84A}{navy}
\fm@curs{2eso}{3B82F6}{white}
\fm@curs{3eso}{FF7A50}{white}
\fm@curs{4eso-a}{A78BFA}{white}
\fm@curs{4eso-b}{6D28D9}{white}
\fm@curs{1batx-mi}{9B1C1C}{white}
\fm@curs{1batx-ccss}{EF6B6B}{white}
\fm@curs{2batx-mii}{1E6B47}{white}
\fm@curs{2batx-ccss}{2BB3A3}{white}

\newcommand{\fmcurs}[1]{%
  \@ifundefined{fm@color@#1}{\PackageError{fentmates}{Curs desconegut: #1}{}}{}%
  \edef\fm@tmp{\noexpand\definecolor{curs}{HTML}{\@nameuse{fm@color@#1}}%
               \noexpand\colorlet{curstext}{\@nameuse{fm@text@#1}}}\fm@tmp
  \colorlet{cursfosc}{curs!72!black}%
  \def\fm@logo{\fm@logodir/vertical-#1.pdf}%
  \def\fm@marca{\fm@logodir/marca-#1.pdf}}

% ---------- capçalera de la primera pàgina ----------
% #1 matèria (petit), #2 curs, #3 tipus de document i unitat, #4 títol
\newcommand{\capcalera}[4]{%
  \gdef\fm@cursnom{#2}%
  \par\noindent
  \begin{tcolorbox}[colback=curs,colframe=curs,boxrule=0pt,arc=2.2mm,width=3.5cm,height=2.15cm,
      valign=center,left=3mm,right=2mm,top=1mm,bottom=1mm,nobeforeafter,box align=center]
    \color{curstext}{\footnotesize #1}\par\vspace{1pt}{\fontsize{19}{21}\selectfont\bfseries #2\par}%
  \end{tcolorbox}\hspace{5mm}%
  \begin{minipage}[c]{\dimexpr\linewidth-3.5cm-5mm-2.9cm\relax}
    {\color{cursfosc}\small\bfseries\MakeUppercase{#3}}\par\vspace{2pt}
    {\color{navy}\fontsize{20}{23}\selectfont\bfseries #4\par}
  \end{minipage}\hfill
  \begin{minipage}[c]{2.6cm}\raggedleft\includegraphics[width=2.6cm]{\fm@logo}\end{minipage}\par
  \vspace{2pt}}

% ---------- barres de secció del color del curs ----------
\renewcommand{\seccio}[2]{\par\needspace{\fmneed\baselineskip}\bigskip
  \begin{tcolorbox}[colback=curs,colframe=curs,coltext=curstext,arc=2mm,boxrule=0pt,left=3mm,right=3mm,top=1.2mm,bottom=1.2mm]
  {\large\bfseries #1\quad #2}
  \end{tcolorbox}\vspace{-2pt}}

% ---------- estrelles de dificultat del color del curs (a Batxillerat no se'n posen) ----------
\@ifundefined{estrelles}{}{%
  \renewcommand{\estrelles}[1]{{\color{curs}\ifcase#1\or\faStar\faStar[regular]\faStar[regular]\or\faStar\faStar\faStar[regular]\or\faStar\faStar\faStar\fi}}}

% ---------- capçalera i peu de les pàgines ----------
% primera pàgina: logos institucionals al peu
\fancypagestyle{primera}{\fancyhf{}\renewcommand{\headrulewidth}{0pt}
  \fancyfoot[L]{\raisebox{-0.32\height}{\includegraphics[height=0.78cm]{logo-gva.png}}\hspace{5mm}%
    \raisebox{-0.32\height}{\includegraphics[height=0.78cm]{logo-vermellar.png}}}
  \fancyfoot[R]{\scriptsize\color{muted}\href{https://creativecommons.org/licenses/by-nc-sa/4.0/deed.ca}{\ccbyncsa}\ \ CC BY-NC-SA 4.0 · fentmates.org\qquad\small\thepage}}
% la resta: el símbol a l'esquerra i el curs a la dreta, amb el títol al centre com abans
\fancyhead[L]{\raisebox{-0.3\height}{\includegraphics[height=0.62cm]{\fm@marca}}}
\fancyhead[R]{\small\bfseries\color{cursfosc}\fm@cursnom}
\renewcommand{\headrule}{\color{curs!45}\hrule height \headrulewidth}
\makeatother

%   \fmcurs{1eso}
% i, al començament del document:
%   \capcalera{Matemàtiques}{1r ESO}{Fitxa d'exercicis · Unitat 2}{Divisibilitat}
% Opcional abans de \input: \def\fmneed{14} (espai mínim abans de cada secció).
% =====================================================================
\makeatletter
\def\fm@logodir{../../logo/pdf}
\providecommand{\fmneed}{10}
\def\fm@cursnom{}

% clau del curs, color i color del text damunt del color
\newcommand{\fm@curs}[3]{\@namedef{fm@color@#1}{#2}\@namedef{fm@text@#1}{#3}}
\fm@curs{1eso}{E6B84A}{navy}
\fm@curs{2eso}{3B82F6}{white}
\fm@curs{3eso}{FF7A50}{white}
\fm@curs{4eso-a}{A78BFA}{white}
\fm@curs{4eso-b}{6D28D9}{white}
\fm@curs{1batx-mi}{9B1C1C}{white}
\fm@curs{1batx-ccss}{EF6B6B}{white}
\fm@curs{2batx-mii}{1E6B47}{white}
\fm@curs{2batx-ccss}{2BB3A3}{white}

\newcommand{\fmcurs}[1]{%
  \@ifundefined{fm@color@#1}{\PackageError{fentmates}{Curs desconegut: #1}{}}{}%
  \edef\fm@tmp{\noexpand\definecolor{curs}{HTML}{\@nameuse{fm@color@#1}}%
               \noexpand\colorlet{curstext}{\@nameuse{fm@text@#1}}}\fm@tmp
  \colorlet{cursfosc}{curs!72!black}%
  \def\fm@logo{\fm@logodir/vertical-#1.pdf}%
  \def\fm@marca{\fm@logodir/marca-#1.pdf}}

% ---------- capçalera de la primera pàgina ----------
% #1 matèria (petit), #2 curs, #3 tipus de document i unitat, #4 títol
\newcommand{\capcalera}[4]{%
  \gdef\fm@cursnom{#2}%
  \par\noindent
  \begin{tcolorbox}[colback=curs,colframe=curs,boxrule=0pt,arc=2.2mm,width=3.5cm,height=2.15cm,
      valign=center,left=3mm,right=2mm,top=1mm,bottom=1mm,nobeforeafter,box align=center]
    \color{curstext}{\footnotesize #1}\par\vspace{1pt}{\fontsize{19}{21}\selectfont\bfseries #2\par}%
  \end{tcolorbox}\hspace{5mm}%
  \begin{minipage}[c]{\dimexpr\linewidth-3.5cm-5mm-2.9cm\relax}
    {\color{cursfosc}\small\bfseries\MakeUppercase{#3}}\par\vspace{2pt}
    {\color{navy}\fontsize{20}{23}\selectfont\bfseries #4\par}
  \end{minipage}\hfill
  \begin{minipage}[c]{2.6cm}\raggedleft\includegraphics[width=2.6cm]{\fm@logo}\end{minipage}\par
  \vspace{2pt}}

% ---------- barres de secció del color del curs ----------
\renewcommand{\seccio}[2]{\par\needspace{\fmneed\baselineskip}\bigskip
  \begin{tcolorbox}[colback=curs,colframe=curs,coltext=curstext,arc=2mm,boxrule=0pt,left=3mm,right=3mm,top=1.2mm,bottom=1.2mm]
  {\large\bfseries #1\quad #2}
  \end{tcolorbox}\vspace{-2pt}}

% ---------- estrelles de dificultat del color del curs (a Batxillerat no se'n posen) ----------
\@ifundefined{estrelles}{}{%
  \renewcommand{\estrelles}[1]{{\color{curs}\ifcase#1\or\faStar\faStar[regular]\faStar[regular]\or\faStar\faStar\faStar[regular]\or\faStar\faStar\faStar\fi}}}

% ---------- capçalera i peu de les pàgines ----------
% primera pàgina: logos institucionals al peu
\fancypagestyle{primera}{\fancyhf{}\renewcommand{\headrulewidth}{0pt}
  \fancyfoot[L]{\raisebox{-0.32\height}{\includegraphics[height=0.78cm]{logo-gva.png}}\hspace{5mm}%
    \raisebox{-0.32\height}{\includegraphics[height=0.78cm]{logo-vermellar.png}}}
  \fancyfoot[R]{\scriptsize\color{muted}\href{https://creativecommons.org/licenses/by-nc-sa/4.0/deed.ca}{\ccbyncsa}\ \ CC BY-NC-SA 4.0 · fentmates.org\qquad\small\thepage}}
% la resta: el símbol a l'esquerra i el curs a la dreta, amb el títol al centre com abans
\fancyhead[L]{\raisebox{-0.3\height}{\includegraphics[height=0.62cm]{\fm@marca}}}
\fancyhead[R]{\small\bfseries\color{cursfosc}\fm@cursnom}
\renewcommand{\headrule}{\color{curs!45}\hrule height \headrulewidth}
\makeatother

%   \fmcurs{1eso}
% i, al començament del document:
%   \capcalera{Matemàtiques}{1r ESO}{Fitxa d'exercicis · Unitat 2}{Divisibilitat}
% Opcional abans de \input: \def\fmneed{14} (espai mínim abans de cada secció).
% =====================================================================
\makeatletter
\def\fm@logodir{../../logo/pdf}
\providecommand{\fmneed}{10}
\def\fm@cursnom{}

% clau del curs, color i color del text damunt del color
\newcommand{\fm@curs}[3]{\@namedef{fm@color@#1}{#2}\@namedef{fm@text@#1}{#3}}
\fm@curs{1eso}{E6B84A}{navy}
\fm@curs{2eso}{3B82F6}{white}
\fm@curs{3eso}{FF7A50}{white}
\fm@curs{4eso-a}{A78BFA}{white}
\fm@curs{4eso-b}{6D28D9}{white}
\fm@curs{1batx-mi}{9B1C1C}{white}
\fm@curs{1batx-ccss}{EF6B6B}{white}
\fm@curs{2batx-mii}{1E6B47}{white}
\fm@curs{2batx-ccss}{2BB3A3}{white}

\newcommand{\fmcurs}[1]{%
  \@ifundefined{fm@color@#1}{\PackageError{fentmates}{Curs desconegut: #1}{}}{}%
  \edef\fm@tmp{\noexpand\definecolor{curs}{HTML}{\@nameuse{fm@color@#1}}%
               \noexpand\colorlet{curstext}{\@nameuse{fm@text@#1}}}\fm@tmp
  \colorlet{cursfosc}{curs!72!black}%
  \def\fm@logo{\fm@logodir/vertical-#1.pdf}%
  \def\fm@marca{\fm@logodir/marca-#1.pdf}}

% ---------- capçalera de la primera pàgina ----------
% #1 matèria (petit), #2 curs, #3 tipus de document i unitat, #4 títol
\newcommand{\capcalera}[4]{%
  \gdef\fm@cursnom{#2}%
  \par\noindent
  \begin{tcolorbox}[colback=curs,colframe=curs,boxrule=0pt,arc=2.2mm,width=3.5cm,height=2.15cm,
      valign=center,left=3mm,right=2mm,top=1mm,bottom=1mm,nobeforeafter,box align=center]
    \color{curstext}{\footnotesize #1}\par\vspace{1pt}{\fontsize{19}{21}\selectfont\bfseries #2\par}%
  \end{tcolorbox}\hspace{5mm}%
  \begin{minipage}[c]{\dimexpr\linewidth-3.5cm-5mm-2.9cm\relax}
    {\color{cursfosc}\small\bfseries\MakeUppercase{#3}}\par\vspace{2pt}
    {\color{navy}\fontsize{20}{23}\selectfont\bfseries #4\par}
  \end{minipage}\hfill
  \begin{minipage}[c]{2.6cm}\raggedleft\includegraphics[width=2.6cm]{\fm@logo}\end{minipage}\par
  \vspace{2pt}}

% ---------- barres de secció del color del curs ----------
\renewcommand{\seccio}[2]{\par\needspace{\fmneed\baselineskip}\bigskip
  \begin{tcolorbox}[colback=curs,colframe=curs,coltext=curstext,arc=2mm,boxrule=0pt,left=3mm,right=3mm,top=1.2mm,bottom=1.2mm]
  {\large\bfseries #1\quad #2}
  \end{tcolorbox}\vspace{-2pt}}

% ---------- estrelles de dificultat del color del curs (a Batxillerat no se'n posen) ----------
\@ifundefined{estrelles}{}{%
  \renewcommand{\estrelles}[1]{{\color{curs}\ifcase#1\or\faStar\faStar[regular]\faStar[regular]\or\faStar\faStar\faStar[regular]\or\faStar\faStar\faStar\fi}}}

% ---------- capçalera i peu de les pàgines ----------
% primera pàgina: logos institucionals al peu
\fancypagestyle{primera}{\fancyhf{}\renewcommand{\headrulewidth}{0pt}
  \fancyfoot[L]{\raisebox{-0.32\height}{\includegraphics[height=0.78cm]{logo-gva.png}}\hspace{5mm}%
    \raisebox{-0.32\height}{\includegraphics[height=0.78cm]{logo-vermellar.png}}}
  \fancyfoot[R]{\scriptsize\color{muted}\href{https://creativecommons.org/licenses/by-nc-sa/4.0/deed.ca}{\ccbyncsa}\ \ CC BY-NC-SA 4.0 · fentmates.org\qquad\small\thepage}}
% la resta: el símbol a l'esquerra i el curs a la dreta, amb el títol al centre com abans
\fancyhead[L]{\raisebox{-0.3\height}{\includegraphics[height=0.62cm]{\fm@marca}}}
\fancyhead[R]{\small\bfseries\color{cursfosc}\fm@cursnom}
\renewcommand{\headrule}{\color{curs!45}\hrule height \headrulewidth}
\makeatother

\fmcurs{1eso}

\begin{document}
\thispagestyle{primera}

\capcalera{Matemàtiques}{1r ESO}{Fitxa d'exercicis · Unitat 10}{Estadística}
\vspace{10pt}
Nom i cognoms: \rule{8.4cm}{0.4pt}\hfill Grup: \rule{1.6cm}{0.4pt}\hfill Data: \rule{2.2cm}{0.4pt}
\vspace{8pt}

\begin{tcolorbox}[colback=soft,colframe=soft,arc=2mm,left=3mm,right=3mm,top=1.5mm,bottom=1.5mm]
\small Exercicis graduats per nivells de dificultat:\quad \estrelles{1}\ bàsic \qquad \estrelles{2}\ mitjà \qquad \estrelles{3}\ avançat.\\[2pt] Resol-los al quadern mostrant tots els passos. Per a fer gràfics, fes servir paper quadriculat, regle i transportador, i posa sempre el títol i el nom de cada eix.
\end{tcolorbox}

% =====================================================
\seccio{1}{Població, mostra i variables estadístiques}

\ex{1} Classifica aquestes variables en qualitatives, quantitatives discretes o quantitatives contínues:\begin{tasks}(3)
  \task nombre de fills
  \task color preferit
  \task temps a córrer 100\,m
  \task pàgines d'un llibre
  \task professió
  \task temperatura d'una ciutat
  \task número de calçat
  \task idioma matern
  \task capacitat d'una botella
\end{tasks}\recorda{Si els valors no són nombres, la variable és \textbf{qualitativa}. Si són nombres: es \textbf{compten} $\to$ \textbf{discreta}; es \textbf{mesuren} $\to$ \textbf{contínua}.}

\ex{1} Per a conéixer quantes hores de televisió veuen a la setmana els 450 alumnes d'un institut, se n'enquesta 60 triats a l'atzar. Indica la població, la mostra, la grandària de la mostra, un individu i la variable estudiada (i de quin tipus és).

\ex{1} En cada cas, digues si s'estudia tota la població o una mostra:\begin{tasks}(1)
  \task La tutora calcula la nota mitjana de la seua classe amb les notes de tots els alumnes.
  \task Per a saber quant duren les piles d'una fàbrica, se'n proven 50.
  \task Una empresa pregunta a 2\,000 persones a qui votaran en les eleccions.
  \task L'entrenador mesura l'alçada de tots els jugadors del seu equip.
\end{tasks}

\ex{2} Explica per què, en aquests estudis, cal treballar amb una mostra i no amb tota la població:\begin{tasks}(1)
  \task la qualitat dels iogurts que fa una fàbrica;
  \task l'opinió de tots els joves d'Espanya sobre les xarxes socials;
  \task el temps que tarden a fondre's les bombetes d'un lot de 10\,000.
\end{tasks}

\ex{2} Raona si aquestes mostres són representatives:\begin{tasks}(1)
  \task Per a saber quantes hores d'esport fan els alumnes d'un institut, es pregunta a 30 persones que ixen del gimnàs.
  \task Per a saber el plat preferit dels alumnes d'un institut, es trien a l'atzar 10 alumnes de cada curs.
  \task Per a saber quina música agrada més a la gent del poble, Laura pregunta als seus amics.
\end{tasks}

\ex{2} Escriu tres valors possibles de cada variable i classifica-la:\begin{tasks}(3)
  \task pes d'un gos
  \task nacionalitat
  \task nombre de cosins
\end{tasks}

\ex{3} Vols saber quant de temps dediquen a llegir cada setmana els 600 alumnes del teu institut. Descriu com faries l'estudi: quina és la població, com triaries una mostra representativa de 60 alumnes, quina pregunta faries i de quin tipus seria la variable.


% =====================================================
\seccio{2}{Freqüències absolutes i relatives}

\ex{1} Aquestes són les mascotes que tenen 20 alumnes: \ $0,\ \allowbreak 1,\ \allowbreak 2,\ \allowbreak 1,\ \allowbreak 0,\ \allowbreak 3,\ \allowbreak 1,\ \allowbreak 1,\ \allowbreak 2,\ \allowbreak 0,\ \allowbreak 1,\ \allowbreak 4,\ \allowbreak 2,\ \allowbreak 1,\ \allowbreak 0,\ \allowbreak 1,\ \allowbreak 2,\ \allowbreak 1,\ \allowbreak 3,\ \allowbreak 1$. \ Fes el recompte i la taula amb les freqüències absolutes, relatives i els percentatges.\recorda{Freqüència absoluta $f_i$: vegades que apareix un valor. Freqüència relativa: $h_i=\dfrac{f_i}{N}$. Percentatge: $h_i\cdot100$.}

\ex{1} Aquests són els grups sanguinis de 25 persones: \ A, O, O, A, B, O, A, AB, O, O, A, B, O, A, O, O, A, B, O, A, AB, O, A, O, A. \ Fes la taula de freqüències absolutes, relatives i percentatges.

\ex{1} Escriu com a percentatge:\begin{tasks}(4)
  \task $0{,}3$
  \task $0{,}05$
  \task $\frac{3}{4}$
  \task $\frac{7}{20}$
\end{tasks}\par I com a freqüència relativa decimal:\begin{tasks}(4)
  \task $12\,\%$
  \task $60\,\%$
  \task $7{,}5\,\%$
  \task $100\,\%$
\end{tasks}

\ex{2} Completa la taula de l'esport preferit de 40 alumnes:\par\smallskip\begin{center}{\renewcommand{\arraystretch}{1.25}\begin{tabular}{|c|C{1.2cm}|C{1.2cm}|C{1.2cm}|}\hline Esport & $f_i$ & $h_i$ & \%\\\hline Futbol & 14 & & \\\hline Bàsquet & & $0{,}25$ & \\\hline Tenis & & & 15\\\hline Natació & 8 & & \\\hline Altres & & & \\\hline \textbf{Total} & 40 & & \\\hline\end{tabular}}\end{center}

\ex{2} Hem llançat un dau 60 vegades. Ha eixit l'1 nou vegades, el 2 onze, el 3 deu, el 4 huit, el 5 dotze i el 6 la resta. Quantes vegades ha eixit el 6? Quin percentatge dels llançaments ha eixit un nombre parell? I un nombre més gran que 4?

\ex{2} En un institut de 120 alumnes, el $35\,\%$ va a peu. Quants alumnes són? Si 18 alumnes van en bici, quina és la freqüència relativa d'aquest valor?

\ex{3} En una classe, la freqüència relativa dels alumnes que porten ulleres és $0{,}2$, i en són 6. Quants alumnes hi ha a la classe? Quants no porten ulleres?

\ex{3} Una variable pren tres valors, $A$, $B$ i $C$. La freqüència relativa d'$A$ és $0{,}25$; la de $B$, $0{,}4$, i la freqüència absoluta de $B$ és 16. Calcula $N$, la freqüència relativa de $C$ i totes les freqüències absolutes.


% =====================================================
\seccio{3}{Taules de freqüències}

\ex{1} Aquestes són les notes d'un examen de 30 alumnes: \ $5,\ \allowbreak 7,\ \allowbreak 6,\ \allowbreak 4,\ \allowbreak 6,\ \allowbreak 8,\ \allowbreak 5,\ \allowbreak 6,\ \allowbreak 3,\ \allowbreak 7,\ \allowbreak 5,\ \allowbreak 9,\ \allowbreak 6,\ \allowbreak 2,\ \allowbreak 7,\ \allowbreak 6,\ \allowbreak 5,\ \allowbreak 8,\ \allowbreak 4,\ \allowbreak 6,\ \allowbreak 10,\ \allowbreak 7,\ \allowbreak 5,\ \allowbreak 3,\ \allowbreak 6,\ \allowbreak 8,\ \allowbreak 4,\ \allowbreak 7,\ \allowbreak 5,\ \allowbreak 9$. \ Fes la taula de freqüències amb $f_i$, $F_i$, $h_i$ i percentatges.\recorda{La freqüència absoluta acumulada $F_i$ és el nombre de dades menors o iguals que $x_i$: se sumen les $f_i$ fins a eixa fila. L'última és $N$.}

\ex{1} Amb la taula de l'exercici anterior, respon:\begin{tasks}(2)
  \task Quants alumnes han tret un 6 o menys?
  \task Quants han aprovat (5 o més)?
  \task Quants han tret més d'un 7?
  \task Quin percentatge ha tret un 9 o més?
\end{tasks}

\ex{2} En aquesta taula, $x_i$ és el nombre de llibres llegits en un mes i $F_i$, la freqüència acumulada. Troba les freqüències absolutes. Quants alumnes han llegit exactament 3 llibres? I almenys 2?\par\smallskip\begin{center}{\renewcommand{\arraystretch}{1.3}\begin{tabular}{|c|C{0.9cm}|C{0.9cm}|C{0.9cm}|C{0.9cm}|C{0.9cm}|C{0.9cm}|}\hline $x_i$ & $0$ & $1$ & $2$ & $3$ & $4$ & $5$\\\hline $F_i$ & $2$ & $5$ & $9$ & $15$ & $18$ & $20$\\\hline\end{tabular}}\end{center}

\ex{2} Aquests són els pesos (en kg) de 24 alumnes: \ $42,\ \allowbreak 51,\ \allowbreak 38,\ \allowbreak 47,\ \allowbreak 55,\ \allowbreak 44,\ \allowbreak 49,\ \allowbreak 36,\ \allowbreak 52,\ \allowbreak 46,\ \allowbreak 41,\ \allowbreak 58,\ \allowbreak 48,\ \allowbreak 45,\ \allowbreak 50,\ \allowbreak 39,\ \allowbreak 53,\ \allowbreak 47,\ \allowbreak 43,\ \allowbreak 56,\ \allowbreak 49,\ \allowbreak 40,\ \allowbreak 45,\ \allowbreak 51$. \ Agrupa'ls en intervals d'amplària 5 des de 35\,kg i fes la taula amb les marques de classe, $f_i$, $F_i$ i percentatges.\recorda{L'interval $[40,\ 45)$ conté les dades majors o iguals que 40 i menors que 45. La marca de classe és el punt mitjà: $\frac{40+45}{2}=42{,}5$.}

\ex{2} En quin interval de l'exercici anterior posaries cadascuna d'aquestes dades?\begin{tasks}(4)
  \task $45$\,kg
  \task $44{,}9$\,kg
  \task $55$\,kg
  \task $35$\,kg
\end{tasks}

\ex{3} Aquests són els temps (en segons) de 20 alumnes en una cursa de 100\,m: \ $13{,}8,\ \allowbreak 15{,}2,\ \allowbreak 14{,}1,\ \allowbreak 12{,}6,\ \allowbreak 16{,}3,\ \allowbreak 14{,}9,\ \allowbreak 13{,}2,\ \allowbreak 15{,}7,\ \allowbreak 14{,}4,\ \allowbreak 13{,}9,\ \allowbreak 15{,}1,\ \allowbreak 12{,}9,\ \allowbreak 14{,}6,\ \allowbreak 16{,}8,\ \allowbreak 13{,}5,\ \allowbreak 14{,}2,\ \allowbreak 15{,}5,\ \allowbreak 14{,}8,\ \allowbreak 13{,}1,\ \allowbreak 15{,}9$. \ Agrupa'ls en intervals d'1 segon des de 12\,s i fes la taula amb $f_i$ i $F_i$. Quin percentatge ha tardat menys de 14\,s? Quants han tardat 15\,s o més?

\ex{3} En una taula de freqüències, $F_3=12$, $F_4=19$ i $F_5=25$, i 5 és el valor més gran. Quantes dades hi ha? Quina és la freqüència absoluta del 4? Quin percentatge de les dades és més gran que 3?


% =====================================================
\seccio{4}{Diagrames de barres, polígons i histogrames}

\ex{1} Dibuixa el diagrama de barres dels grups sanguinis de l'exercici 9 (A: 9, O: 11, B: 3, AB: 2).\recorda{En un \textbf{diagrama de barres} posem els valors en l'eix horitzontal i, damunt de cadascun, una barra d'altura igual a la seua freqüència. Les barres van separades.}

\ex{1} Aquest diagrama mostra quants viatges han fet en avió els alumnes d'una classe. Fes la taula de freqüències. Quants alumnes hi ha? Quina és la moda? Quants han viatjat en avió alguna vegada?\par\smallskip\begin{center}\begin{tikzpicture}[baseline=(current bounding box.center),line cap=round,line join=round]
\draw[muted!30,line width=0.3pt] (0,0.000)--(5.600,0.000); \node[left,font=\tiny,text=muted,inner sep=1.5pt] at (0,0.000) {$0$};
\draw[muted!30,line width=0.3pt] (0,0.350)--(5.600,0.350); \node[left,font=\tiny,text=muted,inner sep=1.5pt] at (0,0.350) {$1$};
\draw[muted!30,line width=0.3pt] (0,0.700)--(5.600,0.700); \node[left,font=\tiny,text=muted,inner sep=1.5pt] at (0,0.700) {$2$};
\draw[muted!30,line width=0.3pt] (0,1.050)--(5.600,1.050); \node[left,font=\tiny,text=muted,inner sep=1.5pt] at (0,1.050) {$3$};
\draw[muted!30,line width=0.3pt] (0,1.400)--(5.600,1.400); \node[left,font=\tiny,text=muted,inner sep=1.5pt] at (0,1.400) {$4$};
\draw[muted!30,line width=0.3pt] (0,1.750)--(5.600,1.750); \node[left,font=\tiny,text=muted,inner sep=1.5pt] at (0,1.750) {$5$};
\draw[muted!30,line width=0.3pt] (0,2.100)--(5.600,2.100); \node[left,font=\tiny,text=muted,inner sep=1.5pt] at (0,2.100) {$6$};
\draw[muted!30,line width=0.3pt] (0,2.450)--(5.600,2.450); \node[left,font=\tiny,text=muted,inner sep=1.5pt] at (0,2.450) {$7$};
\draw[muted!30,line width=0.3pt] (0,2.800)--(5.600,2.800); \node[left,font=\tiny,text=muted,inner sep=1.5pt] at (0,2.800) {$8$};
\draw[-latex,muted,line width=0.6pt] (0,0)--(5.900,0) node[right,font=\scriptsize,text=navy]{Viatges};
\draw[-latex,muted,line width=0.6pt] (0,0)--(0,3.150) node[above,font=\scriptsize,text=navy]{Alumnes};
\draw[teal,fill=teal!35] (0.280,0) rectangle (0.820,1.050);
\node[below,font=\scriptsize,inner sep=2pt] at (0.550,0) {0};
\draw[teal,fill=teal!35] (1.180,0) rectangle (1.720,2.800);
\node[below,font=\scriptsize,inner sep=2pt] at (1.450,0) {1};
\draw[teal,fill=teal!35] (2.080,0) rectangle (2.620,2.100);
\node[below,font=\scriptsize,inner sep=2pt] at (2.350,0) {2};
\draw[teal,fill=teal!35] (2.980,0) rectangle (3.520,1.400);
\node[below,font=\scriptsize,inner sep=2pt] at (3.250,0) {3};
\draw[teal,fill=teal!35] (3.880,0) rectangle (4.420,0.700);
\node[below,font=\scriptsize,inner sep=2pt] at (4.150,0) {4};
\draw[teal,fill=teal!35] (4.780,0) rectangle (5.320,0.350);
\node[below,font=\scriptsize,inner sep=2pt] at (5.050,0) {5};
\end{tikzpicture}\end{center}

\ex{1} Dibuixa el diagrama de barres i el polígon de freqüències del nombre de mascotes de l'exercici 8.

\ex{2} Dibuixa l'histograma i el polígon de freqüències dels pesos agrupats en intervals de l'exercici 19.\recorda{En un \textbf{histograma} les barres ocupen tot l'interval i van juntes. El polígon uneix els punts de dalt sobre les marques de classe i es tanca en l'eix.}

\ex{2} Quin gràfic triaries per a representar cada variable?\begin{tasks}(1)
  \task el nombre de germans de 30 alumnes;
  \task l'alçada de 200 persones;
  \task el mitjà de transport preferit per a anar de vacances;
  \task els minuts que tarden 50 alumnes a fer els deures.
\end{tasks}

\ex{2} Les notes d'un examen de dues classes de 30 alumnes estan en la taula. Dibuixa els dos polígons de freqüències en el mateix gràfic. Quina classe ha tret millors notes?\par\smallskip\begin{center}{\renewcommand{\arraystretch}{1.25}\begin{tabular}{|c|c|c|c|c|c|c|c|}\hline Nota & 4 & 5 & 6 & 7 & 8 & 9 & 10\\\hline Classe A & 1 & 2 & 5 & 7 & 8 & 4 & 3\\\hline Classe B & 0 & 1 & 3 & 6 & 9 & 7 & 4\\\hline\end{tabular}}\end{center}

\ex{3} Tres candidats d'unes eleccions han tret 96, 92 i 98 vots. Un cartell els representa amb un diagrama de barres en què l'eix vertical comença en 90. Per què és enganyós? Dibuixa'l bé.


% =====================================================
\seccio{5}{Diagrames de sectors i pictogrames}

\ex{1} S'ha preguntat a 36 alumnes pel seu esport preferit: futbol 15, bàsquet 9, tenis 6, natació 4 i altres 2. Calcula l'angle de cada sector i dibuixa el diagrama de sectors.\recorda{Angle del sector $=\dfrac{f_i}{N}\cdot360^\circ$. Si $N$ és divisor de 360, cada dada val $360^\circ:N$.}

\ex{1} Calcula els angles del diagrama de sectors dels grups sanguinis de l'exercici 9 (A: 9, O: 11, B: 3, AB: 2) i dibuixa'l.

\ex{2} S'ha preguntat a 200 joves quin aparell fan servir més per a jugar. Quantes persones han triat cada aparell?\par\smallskip\begin{center}\begin{tikzpicture}[baseline=(current bounding box.center)]
\fill[teal!70,draw=white,line width=0.8pt] (0,0)--(90.000:1.45) arc[start angle=90.000,end angle=-54.000,radius=1.45cm]--cycle;
\node[font=\scriptsize\bfseries,text=white] at (18.000:0.899) {$40\,\%$};
\fill[accent!70,draw=white,line width=0.8pt] (0,0)--(-54.000:1.45) arc[start angle=-54.000,end angle=-144.000,radius=1.45cm]--cycle;
\node[font=\scriptsize\bfseries,text=white] at (-99.000:0.899) {$25\,\%$};
\fill[violet!70,draw=white,line width=0.8pt] (0,0)--(-144.000:1.45) arc[start angle=-144.000,end angle=-216.000,radius=1.45cm]--cycle;
\node[font=\scriptsize\bfseries,text=white] at (-180.000:0.899) {$20\,\%$};
\fill[gold!70,draw=white,line width=0.8pt] (0,0)--(-216.000:1.45) arc[start angle=-216.000,end angle=-270.000,radius=1.45cm]--cycle;
\node[font=\scriptsize\bfseries,text=white] at (-243.000:0.899) {$15\,\%$};
\fill[teal!70] (1.85,0.415) rectangle ++(0.22,0.22); \node[right,font=\scriptsize] at (2.10,0.525) {Mòbil};
\fill[accent!70] (1.85,0.065) rectangle ++(0.22,0.22); \node[right,font=\scriptsize] at (2.10,0.175) {Consola};
\fill[violet!70] (1.85,-0.285) rectangle ++(0.22,0.22); \node[right,font=\scriptsize] at (2.10,-0.175) {Ordinador};
\fill[gold!70] (1.85,-0.635) rectangle ++(0.22,0.22); \node[right,font=\scriptsize] at (2.10,-0.525) {Tauleta};
\end{tikzpicture}\end{center}

\ex{2} En un diagrama de sectors sobre el lloc de vacances preferit de 72 persones, els angles són: platja $100^\circ$, muntanya $120^\circ$, ciutat $80^\circ$ i poble $60^\circ$. Quantes persones han triat cada lloc?\recorda{Del sector a la freqüència: $f_i=\dfrac{\text{angle}}{360^\circ}\cdot N$.}

\ex{2} Aquest pictograma mostra les bicis que ha venut una botiga cada mes. Cada icona representa 4 bicis. Quantes n'ha venut cada mes? I en total?\par\smallskip\begin{center}\begin{tikzpicture}[baseline=(current bounding box.center)]
\node[left,font=\small] at (0,-0.00) {Gener};
\node[text=accent,font=\large] at (0.25,-0.00) {\faBicycle};
\node[text=accent,font=\large] at (0.85,-0.00) {\faBicycle};
\node[text=accent,font=\large] at (1.45,-0.00) {\faBicycle};
\node[left,font=\small] at (0,-0.55) {Febrer};
\node[text=accent,font=\large] at (0.25,-0.55) {\faBicycle};
\node[text=accent,font=\large] at (0.85,-0.55) {\faBicycle};
\begin{scope}\clip (1.15,-0.85) rectangle (1.450,-0.25); \node[text=accent,font=\large] at (1.45,-0.55) {\faBicycle};\end{scope}
\node[left,font=\small] at (0,-1.10) {Març};
\node[text=accent,font=\large] at (0.25,-1.10) {\faBicycle};
\node[text=accent,font=\large] at (0.85,-1.10) {\faBicycle};
\node[text=accent,font=\large] at (1.45,-1.10) {\faBicycle};
\node[text=accent,font=\large] at (2.05,-1.10) {\faBicycle};
\node[left,font=\small] at (0,-1.65) {Abril};
\node[text=accent,font=\large] at (0.25,-1.65) {\faBicycle};
\node[text=accent,font=\large] at (0.85,-1.65) {\faBicycle};
\node[text=accent,font=\large] at (1.45,-1.65) {\faBicycle};
\node[text=accent,font=\large] at (2.05,-1.65) {\faBicycle};
\node[text=accent,font=\large] at (2.65,-1.65) {\faBicycle};
\begin{scope}\clip (2.95,-1.95) rectangle (3.250,-1.35); \node[text=accent,font=\large] at (3.25,-1.65) {\faBicycle};\end{scope}
\end{tikzpicture}\end{center}

\ex{2} Els alumnes de quatre grups de 1r d'ESO han recollit taps per a una campanya solidària: 1r A, 20 bosses; 1r B, 35; 1r C, 15, i 1r D, 30. Fes un pictograma en què cada icona valga 5 bosses.

\ex{3} En un diagrama de sectors, el sector de les persones que prefereixen l'estiu fa $72^\circ$ i correspon a 18 persones. Quantes persones s'han enquestat? Si el $30\,\%$ prefereix la primavera, quantes persones són i quin angle té el seu sector?


% =====================================================
\seccio{6}{Moda, mitjana i mediana}

\ex{1} Calcula la moda, la mitjana i la mediana de: \ $5,\ \allowbreak 8,\ \allowbreak 3,\ \allowbreak 8,\ \allowbreak 6,\ \allowbreak 9,\ \allowbreak 7$.\recorda{\textbf{Moda}: el valor més repetit. \textbf{Mitjana}: $\bar x=\dfrac{\text{suma de les dades}}{N}$. \textbf{Mediana}: el valor central de les dades \textbf{ordenades}; si n'hi ha dos, la seua mitjana.}

\ex{1} Calcula la moda, la mitjana i la mediana de: \ $12,\ \allowbreak 15,\ \allowbreak 9,\ \allowbreak 18,\ \allowbreak 15,\ \allowbreak 11$.

\ex{1} S'ha preguntat a 30 alumnes per la seua estació preferida: primavera 9, estiu 13, tardor 3 i hivern 5. Quina és la moda? Es poden calcular la mitjana i la mediana?

\ex{2} Calcula la moda, la mitjana i la mediana de cada conjunt de dades:\begin{tasks}(1)
  \task $4,\ \allowbreak 9,\ \allowbreak 6,\ \allowbreak 9,\ \allowbreak 2$
  \task $1,\ \allowbreak 3,\ \allowbreak 3,\ \allowbreak 7,\ \allowbreak 8,\ \allowbreak 8$
  \task $10,\ \allowbreak 20,\ \allowbreak 30,\ \allowbreak 40$
\end{tasks}

\ex{2} Els gols que ha marcat un equip en els últims 8 partits són: \ $2,\ \allowbreak 1,\ \allowbreak 0,\ \allowbreak 3,\ \allowbreak 2,\ \allowbreak 2,\ \allowbreak 1,\ \allowbreak 4$. \ Calcula la moda, la mitjana i la mediana.

\ex{2} Els preus (en milers d'euros) de les cases venudes en un carrer han sigut: 150, 160, 155, 170, 165 i 900. Calcula la mitjana i la mediana. Quin dels dos paràmetres representa millor el preu d'una casa del carrer? Per què?

\ex{3} Escriu cinc nombres que tinguen la mitjana 6, la mediana 5 i la moda 4.

\ex{3} Les notes de Jan són 4, 6, 7 i 9. Quina és la mitjana? Quina nota ha de traure en el següent examen perquè la mitjana siga 7?


% =====================================================
\seccio{7}{Paràmetres a partir d'una taula de freqüències}

\ex{1} El nombre de germans de 25 alumnes és el de la taula. Calcula la mitjana, la mediana i la moda.\par\smallskip\begin{center}{\renewcommand{\arraystretch}{1.3}\begin{tabular}{|c|C{0.9cm}|C{0.9cm}|C{0.9cm}|C{0.9cm}|C{0.9cm}|}\hline $x_i$ & $0$ & $1$ & $2$ & $3$ & $4$\\\hline $f_i$ & $4$ & $10$ & $6$ & $3$ & $2$\\\hline\end{tabular}}\end{center}\recorda{$\bar x=\dfrac{\sum x_i\cdot f_i}{N}$. Per a la mediana, busca la posició central i mira en quina fila l'$F_i$ hi arriba.}

\ex{1} Calcula la mitjana, la mediana i la moda de les notes de l'exercici 16 a partir de la seua taula de freqüències.

\ex{2} Aquesta taula recull el número de calçat de 28 alumnes. Calcula la mitjana, la mediana i la moda.\par\smallskip\begin{center}{\renewcommand{\arraystretch}{1.3}\begin{tabular}{|c|C{0.9cm}|C{0.9cm}|C{0.9cm}|C{0.9cm}|C{0.9cm}|C{0.9cm}|}\hline Número & $35$ & $36$ & $37$ & $38$ & $39$ & $40$\\\hline $f_i$ & $2$ & $5$ & $7$ & $8$ & $4$ & $2$\\\hline\end{tabular}}\end{center}

\ex{2} Els gols marcats per un equip en 30 partits han sigut: 0 gols en 5 partits, 1 gol en 8, 2 en 9, 3 en 6 i 4 en 2. Calcula la mitjana de gols per partit, la mediana i la moda.

\ex{2} Calcula la mitjana aproximada dels pesos agrupats de l'exercici 19 fent servir les marques de classe.

\ex{3} En aquesta taula falta una freqüència. Troba-la sabent que la mitjana és $1{,}5$.\par\smallskip\begin{center}{\renewcommand{\arraystretch}{1.3}\begin{tabular}{|c|C{0.9cm}|C{0.9cm}|C{0.9cm}|C{0.9cm}|}\hline $x_i$ & $0$ & $1$ & $2$ & $3$\\\hline $f_i$ & $4$ & $a$ & $6$ & $4$\\\hline\end{tabular}}\end{center}

\ex{3} Aquest diagrama de barres mostra la nota (de 1 a 5) que 20 clients han posat a un restaurant. Calcula la mitjana, la mediana i la moda.\par\smallskip\begin{center}\begin{tikzpicture}[baseline=(current bounding box.center),line cap=round,line join=round]
\draw[muted!30,line width=0.3pt] (0,0.000)--(4.700,0.000); \node[left,font=\tiny,text=muted,inner sep=1.5pt] at (0,0.000) {$0$};
\draw[muted!30,line width=0.3pt] (0,0.350)--(4.700,0.350); \node[left,font=\tiny,text=muted,inner sep=1.5pt] at (0,0.350) {$1$};
\draw[muted!30,line width=0.3pt] (0,0.700)--(4.700,0.700); \node[left,font=\tiny,text=muted,inner sep=1.5pt] at (0,0.700) {$2$};
\draw[muted!30,line width=0.3pt] (0,1.050)--(4.700,1.050); \node[left,font=\tiny,text=muted,inner sep=1.5pt] at (0,1.050) {$3$};
\draw[muted!30,line width=0.3pt] (0,1.400)--(4.700,1.400); \node[left,font=\tiny,text=muted,inner sep=1.5pt] at (0,1.400) {$4$};
\draw[muted!30,line width=0.3pt] (0,1.750)--(4.700,1.750); \node[left,font=\tiny,text=muted,inner sep=1.5pt] at (0,1.750) {$5$};
\draw[muted!30,line width=0.3pt] (0,2.100)--(4.700,2.100); \node[left,font=\tiny,text=muted,inner sep=1.5pt] at (0,2.100) {$6$};
\draw[muted!30,line width=0.3pt] (0,2.450)--(4.700,2.450); \node[left,font=\tiny,text=muted,inner sep=1.5pt] at (0,2.450) {$7$};
\draw[muted!30,line width=0.3pt] (0,2.800)--(4.700,2.800); \node[left,font=\tiny,text=muted,inner sep=1.5pt] at (0,2.800) {$8$};
\draw[-latex,muted,line width=0.6pt] (0,0)--(5.000,0) node[right,font=\scriptsize,text=navy]{Nota};
\draw[-latex,muted,line width=0.6pt] (0,0)--(0,3.150) node[above,font=\scriptsize,text=navy]{Clients};
\draw[teal,fill=teal!35] (0.280,0) rectangle (0.820,0.700);
\node[below,font=\scriptsize,inner sep=2pt] at (0.550,0) {1};
\draw[teal,fill=teal!35] (1.180,0) rectangle (1.720,1.750);
\node[below,font=\scriptsize,inner sep=2pt] at (1.450,0) {2};
\draw[teal,fill=teal!35] (2.080,0) rectangle (2.620,2.800);
\node[below,font=\scriptsize,inner sep=2pt] at (2.350,0) {3};
\draw[teal,fill=teal!35] (2.980,0) rectangle (3.520,1.050);
\node[below,font=\scriptsize,inner sep=2pt] at (3.250,0) {4};
\draw[teal,fill=teal!35] (3.880,0) rectangle (4.420,0.700);
\node[below,font=\scriptsize,inner sep=2pt] at (4.150,0) {5};
\end{tikzpicture}\end{center}


% =====================================================
\seccio{8}{Rang i desviació mitjana}

\ex{1} Calcula el rang de cada conjunt de dades:\begin{tasks}(3)
  \task $3,\ \allowbreak 12,\ \allowbreak 7,\ \allowbreak 15,\ \allowbreak 9$
  \task $-4,\ \allowbreak 6,\ \allowbreak 2,\ \allowbreak -1,\ \allowbreak 8$
  \task $2{,}5,\ \allowbreak 4{,}8,\ \allowbreak 1{,}3,\ \allowbreak 3{,}6$
\end{tasks}\recorda{Rang $=x_{\max}-x_{\min}$. \ Desviació mitjana: $DM=\dfrac{\sum|x_i-\bar x|}{N}$ (la mitjana de les distàncies a la mitjana).}

\ex{1} Calcula la desviació mitjana de: \ $3,\ \allowbreak 5,\ \allowbreak 7,\ \allowbreak 9,\ \allowbreak 11$.

\ex{2} Calcula la mitjana, el rang i la desviació mitjana de: \ $2,\ \allowbreak 4,\ \allowbreak 4,\ \allowbreak 6,\ \allowbreak 8,\ \allowbreak 12$.

\ex{2} Les notes de dos alumnes en sis exàmens han sigut: Aina, $5,\ \allowbreak 6,\ \allowbreak 5,\ \allowbreak 6,\ \allowbreak 5,\ \allowbreak 6$; Biel, $2,\ \allowbreak 9,\ \allowbreak 3,\ \allowbreak 8,\ \allowbreak 4,\ \allowbreak 7$. Calcula la mitjana, el rang i la desviació mitjana de cadascun. Qui és més regular?

\ex{2} Les temperatures màximes d'una setmana de cinc dies a dues ciutats han sigut: ciutat X, $12,\ \allowbreak 14,\ \allowbreak 15,\ \allowbreak 13,\ \allowbreak 16$\,°C; ciutat Y, $6,\ \allowbreak 20,\ \allowbreak 10,\ \allowbreak 18,\ \allowbreak 16$\,°C. Compara-les amb la mitjana, el rang i la desviació mitjana.

\ex{3} Calcula la mitjana i la desviació mitjana de les dades d'aquesta taula.\par\smallskip\begin{center}{\renewcommand{\arraystretch}{1.3}\begin{tabular}{|c|C{0.9cm}|C{0.9cm}|C{0.9cm}|C{0.9cm}|C{0.9cm}|}\hline $x_i$ & $1$ & $2$ & $3$ & $4$ & $5$\\\hline $f_i$ & $2$ & $3$ & $5$ & $3$ & $2$\\\hline\end{tabular}}\end{center}

\ex{3} Raona si és vertader o fals:\begin{tasks}(1)
  \task Si el rang és 0, totes les dades són iguals.
  \task Si sumem 5 a totes les dades, la desviació mitjana augmenta 5.
  \task Dos conjunts amb la mateixa mitjana tenen la mateixa desviació mitjana.
  \task La desviació mitjana mai no pot ser negativa.
\end{tasks}


% =====================================================
\seccio{9}{Problemes d'estadística}

\ex{1} Marina ha tret $6{,}5$; $7$ i $5{,}5$ en els tres primers exàmens. Quina nota necessita en el quart per a tindre una mitjana de $6{,}5$?\recorda{Si coneixem la mitjana i el nombre de dades, coneixem la suma: \ $\text{suma}=\bar x\cdot N$.}

\ex{1} La mitjana de quatre nombres és 12. Tres d'ells són 10, 15 i 9. Quin és el quart?

\ex{2} En una classe de 25 alumnes la nota mitjana d'un examen ha sigut $6{,}4$ i en una altra de 15 alumnes, $7{,}2$. Quina és la nota mitjana dels 40 alumnes? Per què no és $6{,}8$?

\ex{2} S'ha preguntat a 20 famílies quants fills tenen: \ $0,\ \allowbreak 1,\ \allowbreak 2,\ \allowbreak 2,\ \allowbreak 1,\ \allowbreak 3,\ \allowbreak 2,\ \allowbreak 1,\ \allowbreak 1,\ \allowbreak 2,\ \allowbreak 0,\ \allowbreak 2,\ \allowbreak 3,\ \allowbreak 1,\ \allowbreak 2,\ \allowbreak 4,\ \allowbreak 2,\ \allowbreak 1,\ \allowbreak 2,\ \allowbreak 1$.\begin{tasks}(1)
  \task Fes la taula de freqüències (amb $f_i$, $F_i$ i $h_i$).
  \task Dibuixa'n el diagrama de barres.
  \task Calcula la moda, la mitjana, la mediana i el rang.
  \task Quin percentatge de famílies té més d'un fill? Quin angle tindria el sector «2 fills» en un diagrama de sectors?
\end{tasks}

\ex{2} L'edat mitjana dels 11 jugadors d'un equip és de 22 anys. Si hi afegim l'entrenador, l'edat mitjana dels 12 és de 24 anys. Quants anys té l'entrenador?

\ex{3} Un professor ha calculat que la mitjana de les notes dels 10 alumnes d'un grup és $6{,}3$, però després s'adona que ha copiat un 9 en lloc d'un 4. Quina és la mitjana correcta? Canvia la mediana si les notes eren 3, 4, 5, 5, 6, 7, 7, 8, 9, 9 (amb el 9 equivocat)?

\ex{3} De dues classes es coneixen aquests resultats d'un examen: classe A, $\bar x=6$, $Me=6$, $DM=0{,}8$; classe B, $\bar x=6$, $Me=7$, $DM=2{,}5$. Raona si és vertader o fals:\begin{tasks}(1)
  \task Les dues classes han tret, de mitjana, la mateixa nota.
  \task En la classe B, almenys la meitat dels alumnes ha tret un 7 o més.
  \task Les notes de la classe A són més semblants entre si.
  \task En la classe A hi ha alumnes amb un 10.
\end{tasks}

\vfill
\begin{center}\color{muted}\small\faCheckCircle[regular]\ \ Ànim! Si et costa algun exercici, repassa els apunts interactius de la unitat.\end{center}

\end{document}
